\section{Peripheral eigenvalues of Schwarz maps}

Maps in this section are complex linear. We use the trace adjoint $T^*$, so
$\tr(T^*(\rho)X)=\tr(\rho T(X))$.

\begin{theorem}[Faithful adjoint stationary state]
    \label{thm:schwarz_faithful_adjoint_state}
    \lean{SchwarzMap.exists_posDef_traceAdjoint_fixedPoint,
        SchwarzMap.exists_posDef_traceAdjoint_stationaryState}
    \leanok
    Let $T$ be positive, unital, and irreducible on $\MN{D}$, with $D\geq1$.
    There is a positive-definite matrix $\rho$ with
    $T^*(\rho)=\rho$ and $\tr(\rho)=1$.
\end{theorem}

\begin{proof}\leanok
    Since $T(\Id)=\Id$, the left Perron theorem at eigenvalue one gives a
    positive-definite fixed point of $T^*$. Its trace is positive, so division
    by its trace gives the required stationary state.
\end{proof}

\begin{theorem}[Scalar fixed points of irreducible unital positive maps]
    \label{thm:schwarz_scalar_fixed_points}
    \lean{SchwarzMap.fixed_eq_scalar_of_irreducible_unital}
    \leanok
    If $T$ is positive, unital, and irreducible, every fixed point is $c\Id$
    for some $c\in\C$.
\end{theorem}

\begin{proof}\leanok
    The ordinary Perron eigenspace at eigenvalue one is spanned by the
    positive-definite eigenvector $\Id$.
\end{proof}

\begin{theorem}[Peripheral Schwarz equality and multiplicative identities]
    \label{thm:schwarz_peripheral_domain}
    \lean{SchwarzMap.schwarz_equality_of_peripheral_eigenvector,
        SchwarzMap.mem_multiplicativeDomain_of_peripheral_eigenvector}
    \leanok
    Let $T$ be a positive Schwarz map and let $\rho>0$ satisfy
    $T^*(\rho)=\rho$. If $T(X)=\gamma X$ and $|\gamma|=1$, then
    $T(X^\dagger X)=T(X^\dagger)T(X)=X^\dagger X$, and for every $Y$,
    $T(YX)=T(Y)T(X)$ and $T(XY)=T(X)T(Y)$.
\end{theorem}

\begin{proof}\leanok
    \uses{thm:abstract_md_characterization}
    The Schwarz defect is positive semidefinite and its trace against
    $\rho$ vanishes by invariance and $|\gamma|=1$. Faithfulness makes the
    defect zero. Apply the same argument to $X^\dagger$, then use the
    multiplicative-domain characterization.
\end{proof}

\begin{theorem}[Peripheral unitaries and phase normalization]
    \label{thm:schwarz_peripheral_unitary}
    \lean{SchwarzMap.exists_peripheral_unitary,
        SchwarzMap.exists_normalized_peripheral_unitary}
    \leanok
    Let $T$ be positive, unital, irreducible, and Schwarz. Every peripheral
    eigenvalue $\gamma$ has a unitary eigenvector $U$. If $\gamma$ is a
    primitive $m$-th root of unity, with $m>0$, the eigenvector can be chosen
    to satisfy $U^m=\Id$.
\end{theorem}

\begin{proof}\leanok
    \uses{thm:schwarz_faithful_adjoint_state,
        thm:schwarz_scalar_fixed_points, thm:schwarz_peripheral_domain,
        thm:schwarz_unitary_powers}
    For a nonzero $\gamma$-eigenvector $X$, the Schwarz equality makes
    $X^\dagger X=c\Id$ with $c>0$. Thus $c^{-1/2}X$ is unitary. Its $m$-th
    power is fixed, hence equals $\alpha\Id$ with $|\alpha|=1$. Choose an
    $m$-th root $\beta$ of $\alpha^{-1}$. Then $|\beta|=1$ and
    $(\beta U)^m=\Id$.
\end{proof}

\begin{theorem}[Powers of peripheral unitaries]
    \label{thm:schwarz_unitary_powers}
    \lean{SchwarzMap.map_powers_of_peripheral_unitary}
    \leanok
    Let $T$ be a positive unital Schwarz map. If $U$ is unitary,
    $T(U)=\gamma U$, and $|\gamma|=1$, then
    $T(U^k)=\gamma^kU^k$ for every $k\in\N$.
\end{theorem}

\begin{proof}\leanok
    \uses{thm:abstract_md_characterization}
    Unitality and unitarity give equality in the Schwarz inequality at $U$.
    Right multiplication by $U$ is therefore preserved. Induction on $k$
    proves the identity, starting from $T(\Id)=\Id$.
\end{proof}

\begin{theorem}[Peripheral product closure]
    \label{thm:schwarz_peripheral_products}
    \lean{SchwarzMap.peripheralEigenvalues_mul_mem}
    \leanok
    The peripheral eigenvalues of a positive unital irreducible Schwarz map
    are closed under multiplication.
\end{theorem}

\begin{proof}\leanok
    \uses{thm:schwarz_faithful_adjoint_state,
        thm:schwarz_peripheral_unitary, thm:schwarz_peripheral_domain}
    Choose unitary eigenvectors $U,V$ for $\mu,\nu$. The multiplicative
    identity gives $T(UV)=\mu\nu UV$. The unitary $UV$ is nonzero and
    $|\mu\nu|=1$.
\end{proof}

\begin{theorem}[Simple peripheral eigenspaces]
    \label{thm:schwarz_peripheral_simple}
    \lean{SchwarzMap.peripheral_eigenspace_finrank_eq_one}
    \leanok
    Every peripheral eigenspace of a positive unital irreducible Schwarz map
    has complex dimension one.
\end{theorem}

\begin{proof}\leanok
    \uses{thm:schwarz_faithful_adjoint_state,
        thm:schwarz_peripheral_unitary, thm:schwarz_peripheral_domain,
        thm:schwarz_scalar_fixed_points}
    Choose a unitary $\gamma$-eigenvector $U$. Every other
    $\gamma$-eigenvector $X$ satisfies $T(XU^\dagger)=XU^\dagger$, hence
    $XU^\dagger=c\Id$ and $X=cU$.
\end{proof}

\begin{theorem}[Cyclicity from finite peripheral product closure]
    \label{thm:peripheral_finite_product_cyclicity}
    \lean{peripheralEigenvalues_card_le_finrank,
        peripheralEigenvalues_eq_range_exp_of_mul_mem}
    \leanok
    Let $T$ act on a finite-dimensional complex vector space $V$, and let
    $S$ be its eigenvalues of modulus one. Then $|S|\leq\dim_{\C}V$.
    If $1\in S$ and $S$ is closed under products, put $m=|S|$ and
    $\gamma=\exp(2\pi i/m)$. Then $m>0$, $\gamma$ is a primitive $m$-th
    root, and $S=\{\gamma^j:0\leq j<m\}$.
\end{theorem}

\begin{proof}\leanok
    \uses{thm:peripheral_finite}
    Choose a nonzero eigenvector for each element of $S$. Distinct
    eigenvalues give linearly independent vectors, proving the bound.
    For $\mu\in S$, multiplication by $\mu$ injects the finite set $S$ into
    itself, hence permutes it. Comparing products gives
    $\prod_{\nu\in S}\nu=\mu^m\prod_{\nu\in S}\nu$, so $\mu^m=1$.
    Thus $S$ is contained in the $m$-th roots of unity. Both sets have
    cardinality $m$, so they coincide.
\end{proof}

\begin{theorem}[Peripheral spectrum of irreducible Schwarz maps]
    \label{thm:wolf_irreducible_schwarz_peripheral}
    \lean{SchwarzMap.wolf_theorem_6_6}
    \leanok
    Let $T:\MN{D}\to\MN{D}$ be positive, unital, irreducible, and Schwarz,
    with $D\geq1$. Write $S=\operatorname{spec}(T)\cap\{z:|z|=1\}$.
    There is an integer $1\leq m\leq D^2$ such that, for
    $\gamma=\exp(2\pi i/m)$:
    \begin{enumerate}
        \item $S=\{\gamma^k:k\in\Z_m\}$.
        \item Every eigenspace at an eigenvalue in $S$ is one-dimensional.
        \item A unitary $U$ satisfies $T(U^k)=\gamma^kU^k$ for every
              $k\in\N$, and $U^m=\Id$.
        \item There are mutually orthogonal projections $(P_k)_{k\in\Z_m}$
              resolving $\Id$. They satisfy $U=\sum_k\gamma^kP_k$ and
              $T(P_{k+1})=P_k$.
    \end{enumerate}
    This is \cite[Theorem~6.6]{Wolf2012Quantum}.
\end{theorem}

\begin{proof}\leanok
    \uses{thm:schwarz_peripheral_products,
        thm:peripheral_finite_product_cyclicity,
        thm:schwarz_peripheral_simple, thm:schwarz_peripheral_unitary,
        thm:schwarz_unitary_powers,
        lem:cyclic_projections_from_peripheral_unitary,
        lem:orthogonal_projection_mul_eq_zero_of_sum_eq_one}
    The identity is a nonzero fixed point, so $1\in S$. Product closure
    and finite cyclicity give the canonical generator and
    $m\leq\dim_{\C}\MN{D}=D^2$. Choose its normalized unitary eigenvector.
    Its power identities give the Fourier spectral projections, whose
    cyclic action is $T(P_{k+1})=P_k$. Orthogonal projections resolving the
    identity are pairwise orthogonal. Simplicity follows from the
    one-dimensional eigenspace theorem.
\end{proof}

\section{Covariance and eigenvector shifts}

\begin{lemma}[Multiplication by a peripheral unitary]
    \label{lem:schwarz_peripheral_unitary_multiplication}
    \lean{SchwarzMap.mem_multiplicativeDomain_of_peripheral_unitary,
        SchwarzMap.map_mul_peripheral_unitary,
        SchwarzMap.map_peripheral_unitary_conjTranspose_mul,
        SchwarzMap.map_conj_peripheral_unitary,
        SchwarzMap.eigenvector_shifts_of_peripheral_unitary,
        SchwarzMap.hasEigenvector_shifts_of_peripheral_unitary}
    \leanok
    Let $T:\MN{D}\to\MN{D}$ be complex linear, positive, unital, and
    Schwarz. Suppose $U$ is unitary, $|\gamma|=1$, and $T(U)=\gamma U$.
    Then $U$ belongs to the multiplicative domain of $T$. For every
    $X\in\MN{D}$,
    \[
        T(U^\dagger X)=\overline\gamma\,U^\dagger T(X),
        \qquad T(XU)=\gamma T(X)U,
        \qquad T(U^\dagger XU)=U^\dagger T(X)U.
    \]
    If $X\ne0$ and $T(X)=\mu X$, then $U^\dagger X$ and $XU$ are nonzero
    eigenvectors at $\overline\gamma\mu$ and $\gamma\mu$, respectively.
\end{lemma}

\begin{proof}\leanok
    \uses{thm:abstract_md_characterization}
    Unitality and $U^\dagger U=UU^\dagger=\Id$ give equality in both
    Schwarz inequalities at $U$, because $|\gamma|=1$. The
    multiplicative-domain characterization gives the right multiplication
    identity. Taking adjoints gives the left identity. Composing the two
    cancels the scalar factors, giving covariance. Substituting
    $T(X)=\mu X$ gives the eigenvector identities; multiplication by a
    unitary is invertible and therefore preserves nonzero vectors.
\end{proof}

\begin{proposition}[Covariance and eigensystem of irreducible maps]
    \label{prop:wolf_schwarz_covariance}
    \lean{SchwarzMap.wolf_proposition_6_7}
    \leanok
    Let $T:\MN{D}\to\MN{D}$ be complex linear, positive, unital,
    irreducible, and Schwarz, with $D\geq1$. Choose $1\leq m\leq D^2$,
    $\gamma=\exp(2\pi i/m)$, and the normalized unitary $U$ from
    Theorem~\ref{thm:wolf_irreducible_schwarz_peripheral}. Thus the
    peripheral spectrum is $\{\gamma^k:k\in\Z_m\}$, $T(U)=\gamma U$,
    and $U^m=\Id$. For every $X\in\MN{D}$,
    \[
        T(U^\dagger XU)=U^\dagger T(X)U.
    \]
    If $T(X)=\mu X$, then
    \begin{equation}
        T(U^\dagger X)=\overline\gamma\mu\,U^\dagger X,
        \qquad T(XU)=\gamma\mu\,XU.
        \label{eq:wolf_schwarz_eigenvector_shifts}
    \end{equation}
    This is \cite[Proposition~6.7 and Equation~(6.38)]{Wolf2012Quantum}.
    % Source: ch06_spectral_properties.tex, lines 803--817.
\end{proposition}

\begin{proof}\leanok
    \uses{thm:wolf_irreducible_schwarz_peripheral,
        lem:schwarz_peripheral_unitary_multiplication}
    Theorem~\ref{thm:wolf_irreducible_schwarz_peripheral} supplies the
    canonical generator and normalized unitary. Apply
    Lemma~\ref{lem:schwarz_peripheral_unitary_multiplication} to this
    unitary to obtain covariance and both eigenvector shifts.
\end{proof}


\section{Peripheral cyclicity for finite Kraus maps}

\begin{theorem}[Peripheral eigenvalues: cyclic structure for Kraus maps]%
    \label{thm:peripheral_cyclic_structure}
    \lean{Kraus.peripheralEigenvalues_eq_range_primitiveRoot}
    \leanok
    \uses{def:kraus_map, def:kraus_adjoint_map, def:unital_kraus,
        def:irreducible_cp, def:peripheral_eigenvalues}
    Let $E$ be an irreducible unital Kraus map on $\MN{D}$ with a
    positive-definite adjoint fixed point (trace-preservation is \emph{not}
    required).
    Then there exist $m \ge 1$ and a primitive $m$-th root of unity
    $\gamma$ such that the peripheral eigenvalues of $E$ are exactly
    $\{1,\gamma,\gamma^2,\ldots,\gamma^{m-1}\}$.
\end{theorem}

\begin{proof}\leanok
    \uses{thm:peripheral_finite_product_cyclicity,
        lem:peripheral_mul_closure}
    The identity supplies the eigenvalue one. Apply finite cyclicity to
    the peripheral set and its product closure.
\end{proof}

\section{Peripheral spectral refinements}%
\label{sec:peripheral_refinements}

The trace-preserving refinements below and the cyclic decomposition and
one-dimensionality results for the peripheral spectrum are generic
finite-Kraus-map statements, with no tensor-network content; they are
relocated here from the periodic-decomposition appendix, whose blocking and
canonical-form arguments cite them across the chapter boundary.

\begin{theorem}[Transfer-map peripheral eigenvalues: cyclic structure]
    \label{thm:transfer_peripheral_cyclic_structure}
    \lean{PeripheralSpectrum.peripheral_eigenvalues_cyclic_structure}
    \leanok
    \uses{def:kraus_map, def:unital_kraus, def:kraus_adjoint_map,
        def:irreducible_cp, def:peripheral_eigenvalues}
    Let $\E_K$ be the Kraus map of a finite matrix family $K$. Assume that
    $\E_K$ is irreducible and unital and has a positive-definite adjoint fixed
    point. Then there exist $m\geq1$ and a primitive $m$-th root of unity
    $\gamma$ such that the peripheral eigenvalues of $\E_K$ are exactly
    $\{1,\gamma,\gamma^2,\ldots,\gamma^{m-1}\}$.
\end{theorem}

\begin{proof}
    \leanok
    \uses{thm:one_peripheral, thm:peripheral_finite,
        lem:peripheral_closed_under_mul_group,
        lem:peripheral_closed_under_inv_group}
    The peripheral eigenvalues form a finite subgroup of $\C^\times$: they
    contain $1$, are nonzero, and are closed under multiplication and
    inversion. Every finite subgroup of $\C^\times$ is cyclic. Choosing a
    generator $\gamma$ gives the asserted primitive root and enumeration.
\end{proof}

\begin{theorem}[Transfer-map peripheral eigenvalues form a cyclic group]%
    \label{thm:peripheral_cyclic_group}
    \lean{PeripheralSpectrum.peripheral_eigenvalues_form_cyclic_group}
    \leanok
    \uses{def:kraus_map, def:unital_kraus, def:tp_kraus,
        def:kraus_adjoint_map, def:irreducible_cp,
        def:peripheral_eigenvalues}
    Let $K$ be a finite matrix family on $\MN{D}$ with Kraus map $\E_K$.
    Assume that $\E_K$ is irreducible, that $K$ is unital and trace-preserving,
    and that the adjoint Kraus map has a positive-definite fixed point. Then
    there exist $m\geq1$ and a primitive $m$-th root of unity
    $\gamma$ such that $m\mid D$ and the peripheral eigenvalues of $\E_K$
    are exactly $\{1,\gamma,\gamma^2,\ldots,\gamma^{m-1}\}$.
\end{theorem}

\begin{proof}
    \leanok
    \uses{thm:transfer_peripheral_cyclic_structure,
        thm:cyclic_decomposition_irreducible_schwarz}
    Apply Theorem~\ref{thm:transfer_peripheral_cyclic_structure} for the cyclic
    structure. Divisibility $m\mid D$ follows from the cyclic
    decomposition: $m$ mutually orthogonal projections summing to $\Id$
    each have equal trace by trace preservation, so
    $m\,\tr(P_0)=D$ with $\tr(P_0)\in\N$.
\end{proof}

\begin{theorem}[Transfer-map period divides bond dimension]%
    \label{thm:channel_period_divides_dim}
    \lean{PeripheralSpectrum.channel_period_divides_dim}
    \leanok
    \uses{def:kraus_map, def:unital_kraus, def:tp_kraus,
        def:kraus_adjoint_map, def:irreducible_cp,
        def:peripheral_eigenvalues}
    Let $K$ be a finite matrix family on $\MN{D}$ with Kraus map $\E_K$.
    Assume that $\E_K$ is irreducible, that $K$ is unital and trace-preserving,
    and that the adjoint Kraus map fixes a positive-definite matrix. Let
    $\gamma$ be a primitive $m$-th root of unity such that the peripheral eigenvalues
    are exactly $\{\gamma^k:k\in\Fin{m}\}$.
    Then $m\mid D$.
\end{theorem}

\begin{proof}
    \leanok
    \uses{thm:cyclic_decomposition_irreducible_schwarz}
    The cyclic decomposition gives orthogonal projections
    $P_0,\ldots,P_{m-1}$ with $\sum_kP_k=\Id$ and
    $E(P_{k+1})=P_k$. Trace preservation gives
    $\tr(P_{k+1})=\tr(P_k)$ for every $k$, hence
    \begin{align}
        D=\tr(\Id)=\sum_{k=0}^{m-1}\tr(P_k)=m\,\tr(P_0).
        \notag
    \end{align}
    Since the trace of an orthogonal projection is an integer, $m\mid D$.
\end{proof}

\subsection{Cyclic decomposition of irreducible finite Kraus maps}

The construction --- normalizing a peripheral eigenvalue to a unitary
eigenvector, forming its discrete Fourier (cyclic) spectral projections,
and the supporting corner-algebra apparatus (restriction to an invariant
corner, corner rank, and the compression isometry
$P\MN{D}P\cong\MN{n}$) --- is given in
Section~\ref{sec:app_spectral_periodicity}.

\begin{theorem}[Cyclic decomposition for irreducible finite Kraus maps]
    \label{thm:cyclic_decomposition_irreducible_schwarz}
    \lean{Kraus.exists_cyclic_decomposition_of_irreducible_schwarz}
    \leanok
    \uses{def:kraus_map, def:unital_kraus, def:kraus_adjoint_map,
        def:irreducible_cp, def:peripheral_eigenvalues}
    Let $K$ be a finite unital matrix family whose Kraus map is irreducible,
    and suppose the adjoint Kraus map fixes a positive-definite matrix. If the
    peripheral spectrum is generated by a primitive $m$-th root $\gamma$, then
    there are a unitary $U$ and orthogonal projections
    $P_0,\ldots,P_{m-1}$ summing to $\Id$ such that
    $E_K(U^k)=\gamma^kU^k$, $U^m=\Id$,
    $U=\sum_{k=0}^{m-1}\gamma^kP_k$, and $E_K(P_{k+1})=P_k$.
\end{theorem}

\begin{proof}
    \leanok
    \uses{lem:normalized_peripheral_unitary, lem:peripheral_unitary_powers,
        lem:cyclic_projections_from_peripheral_unitary}
    Normalize an eigenvector for $\gamma$ to a unitary $U$. Its powers remain
    eigenvectors, and the Fourier projections of $U$ give the stated cyclic
    decomposition.
\end{proof}

\begin{theorem}[Peripheral eigenspaces are one-dimensional]
    \label{thm:peripheral_multiplicity_one}
    \lean{PeripheralSpectrum.peripheral_eigenvalue_multiplicity_one}
    \leanok
    \uses{def:kraus_map, def:unital_kraus, def:kraus_adjoint_map,
        def:irreducible_cp, def:peripheral_eigenvalues}
    Let $K$ be a finite unital matrix family whose Kraus map is irreducible,
    and suppose that the adjoint Kraus map fixes a positive-definite matrix.
    For every peripheral eigenvalue $\gamma$ of the transfer map, the
    $\gamma$-eigenspace is one-dimensional.
\end{theorem}

\begin{proof}\leanok
    \uses{lem:peripheral_unitary_eigenvector,
        thm:ks_equality_of_peripheral_eigenvector_of_fixedPoint,
        thm:right_multiplicative_identity,
        lem:irreducible_unital_fixed_points_scalar}
    Given a peripheral unitary eigenvector $U$ for $\gamma$, the
    multiplicative-domain identity gives
    $E(XU^\dagger)=XU^\dagger$ for any $\gamma$-eigenvector $X$.
    By the scalar fixed-point lemma for irreducible unital maps,
    $XU^\dagger=c\Id$, hence $X=cU$ and the eigenspace is one-dimensional.
\end{proof}
